% Lösungen Einheit 1 – Nr. 13–29
\einheitenkopf[le1][1 Linear oder exponentiell]{Einheit 1 von 5 · Lineares und exponentielles Wachstum unterscheiden}

\erg{13}{a) gleicher Betrag \quad b) gleicher Prozentsatz \quad c) gleicher Betrag \quad d) gleicher Prozentsatz \quad e) gleicher Prozentsatz}
\erg{14}{a) Quotient \quad b) Differenz \quad c) Quotient \quad d) Differenz}
\erg{15}{a) A im Punkt $(0\,|\,0)$ \quad b) B bei 2 \quad c) C bei 3 \quad d) A}
\erg{16}{a) 4; 4; 4 – ja \quad b) 2; 4; 8 – nein \quad c) 15; 15; 15 – ja \quad d) 5; 10; 20 – nein \quad e) 0,6; 0,6; 0,6 – ja \quad f) $-18$; $-18$; $-18$ – ja}
\erg{17}{a) 2; 2; 2 – ja \quad b) 3; 3; 3 – ja \quad c) 2; 1,5; $\approx 1{,}33$ – nein \quad d) 1,25; 1,25; 1,25 – ja \quad e) 0,5; 0,5; 0,5 – ja \quad f) 1,1; $\approx 1{,}09$; $\approx 1{,}08$ – nein}
\erg{18}{a) linear: jedes Jahr kommen 60\,€ dazu \quad b) exponentiell: jeder Wert ist das 1,1-Fache des vorigen \quad c) exponentiell: Faktor 1,5, die Zuwächse 32, 48, 72 werden größer}
\erg{19}{a) linear, jedes Jahr derselbe Betrag \quad b) exponentiell, jedes Jahr derselbe Prozentsatz \quad c) exponentielle Abnahme, Faktor 0,75 \quad d) lineare Abnahme, jede Minute derselbe Betrag \quad e) Sie hat nicht recht: Die Zuwächse 12\,000\,€, 12\,600\,€, 13\,230\,€ werden größer; die Quotienten sind alle 1,05, also exponentielles Wachstum.}
\erg{20}{A Gerade \quad B Kurve \quad C Kurve \quad D Gerade}
\erg{21}{a) A \quad b) B \quad c) E \quad d) D}
\erg{22}{a) G ist eine Gerade; bei gleichem Prozentsatz wächst der Bestand aber immer schneller. \quad b) H beginnt im Ursprung, der Anfangswert ist aber 3000. \quad c) J passt; K ist eine Gerade, die Zahl nimmt aber um denselben Prozentsatz ab, nicht um denselben Betrag; L beginnt bei 8000 statt bei 6000.}
\erg{23}{Punkte $(0\,|\,100)$, $(1\,|\,200)$, $(2\,|\,400)$}
\erg{24}{a) $(2\,|\,225)$: in der Mitte zwischen 200 und 250 \quad b) $(3\,|\,337{,}5)$: drei Viertel des Kästchens zwischen 300 und 350 \quad c) steigende, immer steiler werdende Kurve durch die vier Punkte}
\erg{25}{a) z.\,B. waagerecht 1 Kästchen = 1 Jahr (bis 4 bzw. mehr), senkrecht 1 Kästchen = 200\,€ (bis 2000\,€) \quad b) $(0\,|\,2000)$, $(1\,|\,1600)$, $(2\,|\,1280)$, $(3\,|\,1024)$, $(4\,|\,819{,}20)$ \quad c) fallende, flacher werdende Kurve}
\erg{26}{a) Fehler: gleicher Prozentsatz heißt nicht gleicher Betrag. Die Zuwächse 30; 33; 36,3 werden größer, die Quotienten sind alle 1,1 – das Wachstum ist exponentiell. \quad b) Fehler: Das Guthaben beginnt bei 500\,€, nicht bei 0, und wächst exponentiell. Richtig ist eine gekrümmte Kurve, die bei 500\,€ auf der senkrechten Achse beginnt.}
\erg{27}{a) 1. Jahr: $+100\,€$ (auf 1100\,€), 2. Jahr: $+110\,€$ (auf 1210\,€) – der Zuwachs wächst. \quad b) Der Prozentsatz wird jedes Jahr vom neuen, größeren Wert genommen; deshalb wird der Betrag jedes Jahr größer.}
\erg{28}{a) Nein: nach 52 Wochen wäre sie $25 \cdot 1{,}4^{52}\,$cm hoch, fast 10\,000\,km; das Wachstum hört auf. \quad b) Es gibt nur 800 Schüler; wenn fast alle das Gerücht kennen, kann sich die Zahl nicht mehr verdoppeln.}
\erg{29}{a) A: 1060\,€, 1120\,€, 1180\,€ \quad B: 1050\,€, 1102,50\,€, 1157,63\,€ \quad b) nach 9 Jahren (A: 1540\,€, B: 1551,33\,€) \quad c) 5 Jahre: A (1300\,€ gegen 1276,28\,€); 10 Jahre: B (1628,89\,€ gegen 1600\,€)}
